\section{Background and Related Work}
\label{sec:related}

\subsection{Replication and State-Machine Replication}

The obvious baseline for a single ordered incident journal is state-machine replication. Paxos established the consensus foundation for replicated logs in asynchronous systems with failures \cite{lamport1998parttime}, and Raft made the single-leader log model easier to understand and implement \cite{ongaro2014raft}. More recent systems explore other points in the linearizable-replication design space, including broadcast-based replication in Hermes \cite{katsarakis2020hermes}, predictable performance in HoliPaxos \cite{liang2025holipaxos}, and replicated-key linearizability in LARK \cite{gooding2025lark}. Tango is especially relevant because it builds distributed data structures over a shared log and then materializes application views from that log \cite{balakrishnan2013tango}.

The Rescue OIS incident journal follows the shared-log/materialized-state pattern, but its writer is not chosen by a consensus protocol. The writer is the command vehicle designated by the operational workflow. That choice intentionally trades automatic write failover for an authority model that is explainable to operators and auditable after the incident. \Cref{sec:validation} will compare this design against a Raft baseline in the next phase of the NCA submission work.

\subsection{Local-First and CRDT-Based Replication}

Local-first systems prioritize local availability, user control, and continued operation under poor connectivity \cite{kleppmann2019localfirst}. Much of that work relies on coordination-free replication or CRDT-based synchronization, including efficient delta-based variants \cite{almeida2014delta}. Recent work focuses on preserving application-level invariants and verifying safety properties in replicated settings \cite{haas2023lore,kohler2025consistent,jannes2021owebsync}. Disconnected-operation database replication takes a related multi-primary position: continue local work and reconcile when connectivity returns \cite{mucha2023database}.

Those techniques are valuable for naturally mergeable state, but the incident journal contains ordered, authority-bearing events. CRDTs would still require application-specific semantics for conflicting command decisions, resource assignments, or status transitions. The architecture in this paper therefore adopts offline-first availability for reads and queued submissions, while keeping authoritative incident commits under one command-side sequencer.

\subsection{Edge Computing and Intermittent Connectivity}

Public-safety communications research spans legacy narrowband systems, LTE and 5G public-safety services, ad hoc and vehicular networking, and rapidly deployable mesh-style infrastructure \cite{kumbhar2015legacy,hoyhtya2018critical,wang2023overview}. The most directly relevant field evidence comes from public-safety mesh-radio measurements in forested terrain, where full network connectivity was present only about a third of the observation time and vegetation was the dominant predictor of disconnection \cite{zimbelman2022mesh}. Offline crisis-mapping prototypes such as CRIMP similarly treat vehicles as opportunistic carriers between fragmented field networks \cite{paul2019crimp}.

These works motivate a partition-frequent fault model, but they do not prescribe the application-level consistency model for authority-bearing incident state. The architectural contribution here is above the transport layer: the mesh is treated as intermittent transit, while sequencing, idempotency, promotion, and reconciliation remain explicit protocol concerns.
